%! Regression Lines and Residuals — Unit 2, Lesson 2.3 — Lesson Plan
% GRADUAL-RELEASE plan (warm-up / guided notes and practice / debrief / close and start the
% homework), at the COMPACT budget: the packet is ~6 pages and the lesson finishes in 55 minutes.
% THERE IS NO GROUP ACTIVITY in this course, and NO INDEPENDENT PRACTICE SET in the notes: the
% notes are one Main Ideas / Notes table — I do / we do row by row, then the Guided Practice row
% together. The HOMEWORK is the individual practice, started in class in the last 8 minutes.
% Every box is `skillbox` (breakable) with a \boxguard ahead of it; this course has no
% `fixedskillbox`, no tier boxes, and no exit ticket.
\documentclass[10pt]{article}
\usepackage{apstats-article}
\usepackage{apstats-boxes}

\newcommand{\UnitNumberName}{Unit 2: Exploring Two-Variable Data \quad}
\newcommand{\LessonNumberName}{Lesson 2.3: Regression Lines and Residuals}

\begin{document}
\begin{center}
    {\Large\bfseries \CourseName \\
    {\small\bfseries \UnitNumberName \LessonNumberName}}
\end{center}

\begin{tcolorbox}[colback=sky, colframe=navy, boxrule=0.9pt, arc=2mm,
    left=2mm, right=2mm, top=1.5mm, bottom=1.5mm]
    \textbf{Primary Objective:} Students will use a given regression line to calculate a
    predicted response $\hat{y}$ in context, flag a prediction made by extrapolation as
    unreliable, calculate and interpret a residual as $y - \hat{y}$ with its sign read
    correctly, and read a pre-drawn residual plot to judge whether a linear model is
    appropriate.
    \\[2pt]
    \textbf{CED alignment:} Topics 2.6--2.7 \quad$\cdot$\quad Big Idea: DAT (Data Analysis)
    \quad$\cdot$\quad Skills 2.C, 2.B, 2.A \quad$\cdot$\quad LO DAT-1.D, DAT-1.E, DAT-1.F
    \quad$\cdot$\quad EK DAT-1.D.1--1.D.3, 1.E.1--1.E.2, 1.F.1--1.F.2
    \\[2pt]
    \textbf{Lesson model:} \emph{Gradual release --- I do, we do, you do --- all of it inside
    the guided notes.} There is no group activity. The notes name the vocabulary as they teach
    it and deliver the instruction row by row in one \emph{Main Ideas / Notes} table --- each
    idea a definition to complete and a grid of short problems, the first worked by you and the
    rest by the class; the last row, Guided Practice, is worked together with students holding
    the pen; \textbf{the homework is the individual practice}, and students start it in class
    while you circulate. The whole-class debrief is \emph{spoken}.
    \\[2pt]
    \textbf{Scope:} the line is \emph{given} and only \emph{used} today. Where $a$ and $b$ come
    from, interpreting the slope and intercept in context, and $r^2$ are Lesson~2.4; detailed
    residual-plot departures (curves, fanning) are Lesson~2.5. Every display is pre-drawn.
\end{tcolorbox}

\renewcommand{\arraystretch}{1.4}
\boxguard
\begin{skillbox}[Learning Targets \& Key Understandings]{goldbox}
\begin{tabularx}{\linewidth}{@{}X|X@{}}
\textbf{Learning targets (student language)} & \textbf{Key understandings (the ``why'')} \\[2pt]
\hline
\vspace{-6pt}
\begin{itemize}[leftmargin=1.1em, nosep, after=\vspace{2pt}]
  \item I can plug an $x$-value into a regression line and say what the answer predicts, in
        context.
  \item I can tell a predicted value $\hat{y}$ from an actual value $y$.
  \item I can say when a prediction is extrapolation and why not to trust it.
  \item I can calculate a residual and say whether the line predicted too high or too low.
  \item I can read a residual plot and decide whether a line is a good model.
\end{itemize}
&
\vspace{-6pt}
\begin{itemize}[leftmargin=1.1em, nosep, after=\vspace{2pt}]
  \item A regression model $\hat{y} = a + bx$ uses the explanatory variable to predict the
        response; $\hat{y}$ is a prediction, not an observation (EK DAT-1.D.1--1.D.2).
  \item Extrapolation predicts beyond the interval of $x$-values used to build the line, and
        grows less reliable the farther out it goes (EK DAT-1.D.3).
  \item \textbf{Target misconception:} the residual's \emph{sign}. Students compute
        $\hat{y} - y$ instead of $y - \hat{y}$, or compute it correctly and then read a
        positive residual as ``the line was too high.'' Positive means the actual value is
        \emph{above} the line --- the line \emph{under}predicted (EK DAT-1.E.1).
  \item A residual plot graphs residuals against $x$; apparent randomness around 0 is
        evidence that a linear model is appropriate (EK DAT-1.E.2, 1.F.1--1.F.2).
\end{itemize}
\end{tabularx}
\end{skillbox}

\boxguard
\begin{skillbox}[{Vocabulary, Concepts, \& Theorems --- taught directly in the guided notes}]{greenbox}
\begin{minipage}[t]{0.48\linewidth}
    \begin{itemize}
        \item \textbf{Regression line} --- $\hat{y} = a + bx$; uses $x$ to predict $y$
              (EK DAT-1.D.1)
        \item \textbf{Predicted value} $\hat{y}$ --- the line's value at a given $x$
              (EK DAT-1.D.2)
        \item \textbf{Extrapolation} --- predicting outside the $x$-interval of the data;
              less reliable the farther out (EK DAT-1.D.3)
    \end{itemize}
\end{minipage}
\hfill
\begin{minipage}[t]{0.48\linewidth}
    \begin{itemize}
        \item \textbf{Residual} --- $y - \hat{y}$, actual minus predicted (EK DAT-1.E.1)
        \item \textbf{Residual plot} --- residuals vs.\ $x$; random scatter around 0 supports
              a linear model (EK DAT-1.E.2, 1.F.1)
        \item \textbf{Carried forward} --- explanatory/response, direction, form, and $r$
              (Lesson~2.2)
    \end{itemize}
\end{minipage}
\end{skillbox}

% A tabularx never splits, so guard generously ahead of it.
\boxguard[30]
\begin{skillbox}[Lesson at a Glance --- \MeetingLength]{sky}
\renewcommand{\arraystretch}{1.25}
\begin{tabularx}{\linewidth}{@{}l c X X@{}}
\textbf{Phase} & \textbf{Min} & \textbf{Students} & \textbf{Teacher} \\
\hline
Warm-Up & 5 & Describe the pizza scatterplot, plug into $y = 8 + 3x$, and place two points above or below the line & Debrief item~3 aloud; leave ``$(5,26)$ is 3 \emph{above}'' on the board all period \\
Guided Notes \& Practice & 34 & Fill the vocabulary box and the notes table: problems 1--10 row by row, then \emph{Sleep \& Reaction} (11--14) together, pen in their hand & \textbf{I do} each row's definition lines and first problem (1, 4, 7), \textbf{we do} the rest of its grid (24) $\to$ \textbf{we do} Guided Practice (10) \\
Debrief & 8 & Pens down, papers up --- answer aloud, nothing to fill in & Open on an \emph{almost}-right problem~9, then the sign in context from problem~12 \\
Close \& Assign & 8 & \textbf{Start the homework in class}, alone and silent & Announce the paper homework, circulate for the second formative read (A2), preview Lesson~2.4 \\
\end{tabularx}
\end{skillbox}

\boxguard[20]
\begin{skillbox}[Warm-Up --- Activate Prior Knowledge (5 min)]{sky}
\begin{minipage}[t]{0.48\linewidth}
    \textbf{The three items and what each seeds}
    \begin{itemize}
        \item \textbf{1.} The pizza scatterplot in words: explanatory variable, direction,
              form, $r$ near 0.95. \textbf{Spirals} to Lesson 2.2 and loads the context the
              notes open on.
        \item \textbf{2.} Plug into $y = 8 + 3x$. \textbf{Seeds} row~1: prediction is
              substitution --- nothing new, only a new name.
        \item \textbf{3.} Is $(5, 26)$ above or below the line, and by how much?
              \textbf{Seeds} row~3: the residual, before it has a name or a sign convention.
    \end{itemize}
\end{minipage}
\hfill
\begin{minipage}[t]{0.48\linewidth}
    \textbf{Running it}
    \begin{itemize}
        \item \textbf{Debrief item~3 aloud} and write ``$(5,26)$ is 3 \emph{above} the line''
              on the board. Leave it up: problem~7 turns it into ``residual $= +3$,'' and
              problem~9 leans on the word \emph{above}.
        \item Items~1 and~2 are a minute each. Do not let item~1 become a correlation review
              --- one sentence and move on.
        \item Expect some students to answer item~3 with a subtraction in the wrong order
              (``$23 - 26 = -3$, below''). Do not correct it yet --- the picture says above,
              and the crux settles it.
    \end{itemize}
\end{minipage}
\end{skillbox}

\boxguard
\begin{skillbox}[Guided Notes \& Practice --- I do, then we do (34 min)]{sky}
\textbf{This is the whole lesson.} There is no group activity, and there is no independent
practice set on this page: \textbf{the homework is the individual practice}, started in class.
The notes are one \emph{Main Ideas / Notes} table --- four rows, 14 short problems, two pages.
\textbf{In every row you fill the definition lines and work the first problem of the grid
(problems 1, 4, 7); the class works the rest, pen in their hand.} The vocabulary box is filled
\emph{as you go}, never in advance.
\begin{multicols}{2}
\textbf{The Equation \& Prediction (6 min).} Fill the five blanks, write
$\widehat{\text{time}} = 8 + 3(\text{distance})$, and work problem~1 on the scatterplot:
find 5 miles, go up to the line, read 23. Problem~2 is theirs. \textbf{Problem~3 is the trap}:
``A 6-mile delivery took 24 minutes. Give $y$ and $\hat{y}$.'' The 24 is \emph{given}, so half
the room writes $\hat{y} = 24$. Take a hand count before anyone speaks: \emph{``Which number came
from the line?''}

\textbf{Predicting Beyond the Data (6 min).} Problem~4 yourself: 30 miles gives 98 minutes, and
the equation never objects. That is the point --- the line always answers; the data decide
whether to believe it. Problems~5--6 are theirs. Hold them to the reason in problem~6: not ``it
might be wrong'' but ``we have no deliveries out there to check.''

\textbf{Actual Minus Predicted (12 min) --- the must-land row, and the crux.} Fill
residual $= y - \hat{y}$ and point at the warm-up sentence still on the board: $(5,26)$ was 3
\emph{above}, and $26 - 23 = +3$. Work problem~7 that way. Problem~8 is theirs --- a negative
residual, the line \emph{over}predicted. \textbf{Problem~9 is the crux}: Jordan subtracts in
the wrong order \emph{and} reads the sign backwards. Give it two full minutes, silent, and
demand \emph{both} fixes. Many students fix the arithmetic and keep ``too high'' --- that is
the misconception surviving, and it is the paper you want for the debrief. Problem~10 names the
residual plot's one job today: random scatter, no pattern, a line is fine.

\columnbreak
\textbf{We do (about 10 min): Sleep \& Reaction.} Problems 11--14, a fresh context, worked
entirely together. Students hold the pen; you hold the questions: \emph{``Is 270 above or below
255? So is the line too high or too low?''} \quad \emph{``Is 12 hours inside 4 to 9?''}
\textbf{Problem~12 is the crux transferred}, and the context adds one layer: a positive
residual here is a \emph{slower} reaction, so ``positive'' does not mean ``good.'' Protect it
when time is short; problem~13 can be done aloud.

\textbf{The pen must actually change hands.} Put problem~11 on them cold, then take a hand
count on problem~12's sign before anyone explains it.

Circulate and demand the reason, not the number:
\begin{itemize}[leftmargin=1.1em, nosep]
  \item \textbf{Problem~12:} ``Which came first in your subtraction --- the 270 or the 255?
        Why that one?''
  \item \textbf{Problem~13:} ``The equation gave you a number. Why don't you believe it?''
  \item \textbf{Problem~14:} ``What would the plot have to look like for you to say no?''
\end{itemize}
While circulating, pick the papers to put up in the debrief --- one problem~9 with the
arithmetic fixed but ``too high'' kept (the \emph{almost}-right answer), one with both fixed.
\textbf{Your formative read comes twice}: here, and again in the last eight minutes as students
start the homework.
\end{multicols}
\end{skillbox}

\boxguard
\begin{skillbox}[Debrief --- whole class, spoken (8 min)]{sky}
Pens down, papers up. \textbf{This debrief is spoken} --- there is no box to fill in, so it lives
on the board and in the room.
\begin{multicols}{2}
\begin{enumerate}[leftmargin=1.3em, nosep, itemsep=3pt]
  \item \textbf{Open on the \emph{almost}-right problem~9}: $26 - 23 = +3$, ``so the line
        predicted too high.'' Ask the room what is right about it and what is still wrong.
        Then point at the scatterplot: the point is above the line, so the line was
        \emph{under} it --- too low.
  \item Say the rule as one sentence and leave it up: \emph{residual = actual minus predicted;
        positive means the line guessed too low.}
  \item Problem~12: the same sign, a new meaning in context --- 15~ms \emph{slower} than
        predicted. The sign is about the line, not about good or bad.
  \item Problem~13: the equation gave 130~ms for 12 hours. Nothing in the arithmetic warned
        us; only the data range did.
\end{enumerate}
\columnbreak
\textbf{Two things to demand aloud.} First, \textbf{make a student interpret one residual
completely in context} --- the number, the direction, and what the line did. Second,
\textbf{ask what a residual plot would have to look like to make them distrust the line}; take
any ``a pattern, a curve'' answer and tell them Lesson~2.5 is exactly that.

\textbf{If the debrief runs short}, cut item~4 --- item~1 is the one that cannot be cut.
\textbf{Do not let the debrief eat the last eight minutes} --- the homework start is not optional
padding, it is your second formative read.
\end{multicols}
\end{skillbox}

\boxguard
\begin{skillbox}[AP Practice --- assigned, not scored]{goldbox}
One page on a coffee shop's hot-drink sales against daily high temperature,
$\widehat{\text{drinks}} = 210 - 2.4(\text{temperature})$, 30--75$^\circ$F. \textbf{MC~1}: the
residual on a 50$^\circ$F day with 100 drinks sold --- (B) $+10$, sold more than predicted;
distractor (A) is the sign flip and (C) is the right number with the wrong reading.
\textbf{MC~2}: a $-30$-drink prediction at 100$^\circ$F --- (C) extrapolation. \textbf{Free
response}: (a) predict at 60$^\circ$F (66 drinks); (b) residual for 58 drinks sold ($-8$, the
line overpredicted), interpreted in context; \textbf{(c) is the spiral} to Lesson~2.2 ---
negative direction, and association is not causation.

\textbf{Not scored --- the cover's score column reads \textbf{NA} for this page.} Go over it at
the start of Lesson~2.4 and hold the AP standard out loud: \emph{in context or it does not
count}. Part (b) without the words ``fewer hot drinks than predicted'' is incomplete, and part
(c) that says ``yes, heat causes it'' does not earn the point.
\end{skillbox}

\boxguard
\begin{skillbox}[Homework --- scored, due the first class after two study halls]{goldbox}
Five items on used-car prices, $\widehat{\text{price}} = 32.0 - 0.24(\text{mileage})$, 8--95
thousand miles --- the third context. \textbf{Part A}: A1 predict at 50 thousand miles
(\$20{,}000); \textbf{A2} the residual for a car listed at \$22{,}500 ($+2.5$ thousand --- the
line underpredicted); A3 a 150-thousand-mile prediction that comes out negative, because it is
extrapolation. \textbf{Part B}: B1 read a pre-drawn residual plot (random scatter --- a line is
appropriate); B2 an AP-style multiple-choice item on a residual of $-1.8$ with a required
justification.

\textbf{Start it in the last eight minutes of the period, in class and alone.}
\textbf{Scoring:} 10 points --- 2 per item. \textbf{A2 is where the misconception shows}: a
student who writes $-2.5$ or ``the line guessed too high'' has not yet owned the sign.
\textbf{B2 carries the check:} distractors (A) and (C) are both the flipped sign in AP clothing.

\textbf{Paper homework is the default for this course} --- assign this page unless you have
decided otherwise. \textbf{DeltaMath override (occasional):} on the days you assign DeltaMath
instead, it replaces this page, you strike through the score cell on the cover, and this page
stays in the packet as extra practice. Never assign both.
\end{skillbox}

\boxguard
\begin{skillbox}[Watch For (while circulating)]{redbox}
\begin{itemize}[leftmargin=1.2em, nosep]
  \item \textbf{The residual's sign, flipped or misread} (notes~7--9, 12; AP MC~1, (b); HW~A2,
        B2 --- the target misconception). Computing $\hat{y} - y$, or a positive residual read
        as ``the line was too high.'' Probe: \emph{``Is the point above or below the line?
        So is the line above or below the point?''}
  \item \textbf{$y$ and $\hat{y}$ confused} (notes~3; HW~A2): the given data value written as
        the prediction. Probe: \emph{``Which number came out of the equation?''}
  \item \textbf{Extrapolation trusted because the arithmetic worked} (notes~4, 13; AP MC~2;
        HW~A3). Probe: \emph{``What $x$-values was this line built from?''}
  \item \textbf{A residual plot read as a scatterplot} (notes~10, 14; HW~B1): ``the points go
        up, so it's linear.'' Probe: \emph{``What is on the vertical axis here --- and where
        is the zero line?''}
  \item \textbf{Slope interpretation creeping in early} --- welcome the question, park it for
        Lesson~2.4.
\end{itemize}
Cold-call prompts: \emph{``Give me a residual of zero in words.''} \quad \emph{``Name an
$x$-value that would be extrapolation for the pizza line.''}
\end{skillbox}

\boxguard
\begin{skillbox}[Close \& Start the Homework (8 min)]{goldbox}
\textbf{Students start the homework here, in class, alone and silent --- this is the individual
practice, and the first minutes of it are your best formative read.} Hand the launch first ---
five items on paper, scored, due the first class after two study halls --- and point out that
the AP Practice page before it is theirs to work and bring to review. Circulate:
\textbf{A2 is the item that tells you whether the crux landed}. Sort what you see into three
piles as you walk --- $22.5 - 20.0 = +2.5$ and ``listed above the prediction; the line
underpredicted'' (ready); $+2.5$ with no reading, or a reading with no number (ready, needs the
language); $-2.5$ or ``too high'' (the misconception survived) --- and open Lesson~2.4 by
reworking problem~9 if that third pile ran past a third of the room. Then close on what changed:
\emph{``You walked in able to plug a number into a line. You walk out knowing the line is a
guess, how far each point missed it, and which way.''} \textbf{Preview:} Lesson~2.4,
\emph{Least-Squares Regression} --- where $a$ and $b$ come from, what the slope and intercept
say in context, and $r^2$.
\end{skillbox}

% ── Teacher notes — one per component, in packet order (four: there is no activity) ──
\begin{teachernote}[Warm-Up]
Three items, five minutes. Item~3 is the seed: ``is $(5, 26)$ above or below $y = 8 + 3x$, and by
how much?'' is a residual with its sign read from the picture before anyone has a formula to
misapply. \textbf{Write the room's answer --- ``3 above'' --- on the board and leave it all
period}; problems~7 and~9 lean on it. If a student subtracts $23 - 26$ and says ``below,'' ask
the room to look at the height, not the subtraction, and leave the order question for the notes.
Item~1 is a one-minute spiral to Lesson~2.2 --- do not let it grow.
\end{teachernote}

\begin{teachernote}[Guided Notes \& Practice]
\textbf{Pace it 24 / 10, then 8 for the debrief, then 8 for the homework start.} Four rows, 14
problems, two pages. \textbf{You work problems 1, 4, and~7}; the class works the rest of each
grid; problems 11--14 are Guided Practice, together. The first two rows are machinery and must
move --- six minutes each. Prediction is substitution, and students already know how to
substitute; the only new thing in row~1 is the hat, which is why \textbf{problem~3 is the trap}
(the given 24 is $y$, the line's 26 is $\hat{y}$).

\textbf{Spend the time in Actual Minus Predicted}, which carries the misconception.
\textbf{Problem~9 is the crux} --- Jordan's $23 - 26 = -3$, ``so the line predicted too high,''
has two errors, and the second survives the first: a student who fixes the order often keeps
``too high.'' Do not rescue it. Collect one such paper for the debrief.

\textbf{Keep the residual plot modest.} Problem~10 and problem~14 ask one thing: is there a
pattern? Random scatter means a line is appropriate. Curves, fans, and outliers in a residual
plot are Lesson~2.5; if a student asks, say so and move on. Likewise, a student who asks what
the 3 in $8 + 3x$ \emph{means} is asking the first question of Lesson~2.4 --- park it.

\textbf{Sleep \& Reaction is where the pen changes hands.} Problem~12 is the crux transferred,
and the context makes it harder: a positive residual is a \emph{slower} reaction. If you only
get through 11, 12, and~14, that is the right trade; 13 can be done aloud.

\textbf{There is no independent practice set on this page, and that is deliberate --- the
homework is the individual practice.} The period ends with students starting it in front of
you; A2 is the item to watch. \textbf{Debrief (8 min), spoken}: open on the almost-right
problem~9. \textbf{Do not let the debrief eat the last eight minutes.}
\end{teachernote}

\begin{teachernote}[AP Practice]
\textbf{This page is not required and not scored} --- the graded practice for this lesson is
the homework behind it. Work it as AP-format reps and go over it at the start of Lesson~2.4.
\textbf{MC~1 is the one worth reading closely}: (A) is $\hat{y} - y$ and (C) has the right
number with the wrong reading, which is exactly problem~9's two errors split into two
distractors. If more than a handful choose (C), the sign-reading half of the misconception is
the one that survived --- reopen the ``above the line means the line was too low'' picture
before Lesson~2.4 builds residuals into least squares. Part~(c) is the spiral; it must hedge
(association, not causation) to earn credit.
\end{teachernote}

\begin{teachernote}[Homework]
\textbf{Scored, 10 points (2 per item), due the first class after two study halls.} The eight
minutes at the end of the period are a \emph{start}, not a completion: put students on A1, A2,
and~B2 while you can still catch a wrong start, and let A3 and B1 go home. \textbf{A2 is the
one to read first when scoring} --- full credit needs $+2.5$ (thousand dollars), the words
``listed above the prediction,'' and ``underpredicted'' or equivalent; a correct number with
``too high'' earns 1 of 2. A3 earns full credit only if it names extrapolation (150 is outside
8--95), not merely ``prices can't be negative.'' \textbf{B2}: (B), and the justification must
use the sign. If more than a third of the room misses A2 or B2, open Lesson~2.4 with problem~9
again --- Lesson~2.4 minimizes squared residuals, and it rests on this.

\textbf{Paper is the default.} On the occasions you assign DeltaMath in its place, strike the
score cell on the cover and tell students this page is now optional practice. Do not assign
both.
\end{teachernote}
\end{document}
